\documentclass{article}
\usepackage{/globals/math}
\usepackage{/globals/compsci}
\usepackage{/globals/anki}
\ankiprefix{text}
\ankitopic{Text}
\ankishow
\begin{document}
\section{Texte} 
Auch Texte, Charaktere, werden im Computer binär dargestellt. Hier wird ein Standard festgelegt, welche Bitsequenz welchen Charakter darstellt.

Die zwei relevantesten dieser Standards sind ASCII, mit immer 8 Bits pro Zeichen, aber auch nur 128 Zeichen, und Unicode, mit mehreren zehntausenden an Zeichen, aber mit bis zu 4 Bytes pro Zeichen.
\ankinote{encs}{Encodings}{ASCII (1 Byte), Unicode (bis 4 Bytes)}

Dass ein Text, also eine Zeichensequenz, zuende ist wird durch einen 00\bh-Byte gezeigt.
\ankinote{ende}{Ende}{00h-Byte}
\end{document}